\label{ch:nsi-observables}

\section{Named NSI presets}
\label{sec:presets-table}

\modulehead{config/presets.py}{\code{NSI\_PRESETS}: named, literature-justified
$\epsilon$ configurations, built via \pyfunc{NSIConfig.from\_preset}.}

\begin{implbox}
Every preset below sets a handful of $\epsilon$ parameters to a
representative, literature-motivated value (all other components default to
zero) rather than to a unique global best fit: current data leaves large NSI
degeneracies, so no single preferred point exists. Their purpose is
threefold: (1) give the validation notebooks and tests physically motivated
$\epsilon$ values instead of arbitrary numbers, covering diagonal-only,
off-diagonal, single- and multi-parameter regimes; (2) reproduce, and stay
traceable to, specific published sensitivity/bound studies; and (3) let
users reach a standard benchmark scenario (e.g.\ the LMA-Dark degeneracy) by
name instead of re-deriving its $\epsilon$ values. \textbf{None of the nine
presets currently sets any $\epsilon^n$ (neutron-normalized) component}: every
one is, in the language of \cref{sec:theory-composition}, a single-composition
fit with an implicit $\epsilon^n=0$ hypothesis. A user who wants to explore
composition dependence explicitly builds an \pyclass{NSIConfig} with both
blocks set directly, e.g.\ \code{NSIConfig(eps\_ee=0.30, eps\_ee\_n=0.10)}, or
starts from a preset and overrides it with \code{dataclasses.replace} (safe
for a preset-built config, since it was constructed from the scalar fields,
not \code{from\_raw\_epsilon} -- \cref{sec:impl-nsiconfig}).

\begin{center}
\small
\begin{tabular}{@{}ll@{}}
\toprule
Preset & $\epsilon$ values and origin \\
\midrule
\parbox[t]{3.2cm}{\code{sm\_\allowbreak{}no\_\allowbreak{}nsi}} &
  \parbox[t]{9.4cm}{All $\epsilon=0$. Standard Model limit; equivalent to
  passing \code{nsi=None}.} \\[10pt]
\parbox[t]{3.2cm}{\code{nsi\_\allowbreak{}diagonal\_\allowbreak{}biggio2009}} &
  \parbox[t]{9.4cm}{$\epsilon_{ee}=0.30$, $\epsilon_{\tau\tau}=0.15$.
  Illustrative diagonal-only benchmark below the one-parameter limits
  collected by \textcite{BiggioBlennowFM2009}; those limits depend on the
  operator, fermion and normalization conventions.} \\[18pt]
\parbox[t]{3.2cm}{\code{nsi\_\allowbreak{}lma\_\allowbreak{}dark\_\allowbreak{}esteban2018}} &
  \parbox[t]{9.4cm}{$\epsilon_{ee}=-2.0$. The LMA-Dark degenerate solar
  solution of \textcite{Esteban2018LMADark} (\S\ref{sec:theory-lma-dark}):
  this large negative $\epsilon_{ee}$ flips the sign of the effective MSW
  potential, mimicking $\theta_{12}\to\pi/2-\theta_{12}\approx56.6^\circ$
  (paired with the \code{"\_LMA\_DARK\_NUFIT52\_NO"} oscillation preset).
  Validated numerically in \cref{sec:results-lma-dark}.} \\[26pt]
\parbox[t]{3.2cm}{\code{nsi\_\allowbreak{}offdiag\_\allowbreak{}icecube2022}} &
  \parbox[t]{9.4cm}{$\epsilon_{\mu\tau}=0.005$, $\epsilon_{\tau\tau}=0.015$.
  Illustrates the sector and parameter scale probed by the IceCube DeepCore
  atmospheric analysis \parencite{IceCube2022NSI}; this exact pair is not a
  published best fit or confidence-limit point.} \\[18pt]
\parbox[t]{3.2cm}{\code{nsi\_\allowbreak{}dune\_\allowbreak{}etau}} &
  \parbox[t]{9.4cm}{$\epsilon_{e\tau}=0.15$ (real). Illustrative DUNE
  propagation-NSI point motivated by the octant-degeneracy study of
  \textcite{Agarwalla2016DUNENSI}; it is not a DUNE best-fit point.} \\[18pt]
\parbox[t]{3.2cm}{\code{nsi\_\allowbreak{}hyperk\_\allowbreak{}etau}} &
  \parbox[t]{9.4cm}{$\epsilon_{ee}=0.20$, $\epsilon_{e\tau}=0.10$ (real).
  Illustrative combination of a diagonal MSW correction and a
  flavour-changing coupling, motivated by Hyper-Kamiokande NSI-sensitivity
  studies \parencite{Kelly2017HyperKNSI,Fukasawa2017HyperKNSI}; the exact
  pair is not a published best fit.} \\[18pt]
\parbox[t]{3.2cm}{\code{nsi\_\allowbreak{}globalfit\_\allowbreak{}esteban2018}} &
  \parbox[t]{9.4cm}{$\epsilon_{ee}=0.25$, $\epsilon_{e\tau}=0.12$,
  $\epsilon_{\mu\tau}=0.01$, $\epsilon_{\tau\tau}=0.08$. Multi-parameter
  illustrative stress test inspired by the global oscillation analysis of
  \textcite{Esteban2018LMADark}; it is not a published fit point or an
  allowed-region claim, since correlations, degeneracies, and matter
  composition all matter (\S\ref{sec:theory-composition}).} \\[24pt]
\parbox[t]{3.2cm}{\code{nsi\_\allowbreak{}ee\_\allowbreak{}only}} &
  \parbox[t]{9.4cm}{$\epsilon_{ee}=0.30$. Simplest possible NSI scenario:
  only the effective MSW potential is modified, with no flavour-changing
  coupling \parencite{BiggioBlennowFM2009}.} \\[10pt]
\parbox[t]{3.2cm}{\code{nsi\_\allowbreak{}mutau\_\allowbreak{}only}} &
  \parbox[t]{9.4cm}{$\epsilon_{\mu\tau}=0.01$ (real). Illustrative
  atmospheric stress-test just outside the IceCube DeepCore 90\% CL interval
  $-0.0067<\epsilon_{\mu\tau}<0.0081$
  \parencite{IceCube2018MutauNSI}; it is not an allowed-point claim.} \\[14pt]
\bottomrule
\end{tabular}
\end{center}
\end{implbox}

\begin{refbox}
\cite{BiggioBlennowFM2009,Esteban2018LMADark,IceCube2022NSI,Agarwalla2016DUNENSI,Kelly2017HyperKNSI,Fukasawa2017HyperKNSI,IceCube2018MutauNSI}.
\end{refbox}

\section{Building an oscillation scenario with NSI}

\begin{examplebox}[title={Attaching an NSI preset to a propagation config}]
\begin{lstlisting}[style=tpython]
from tpeanuts.config.propagation import PropagationConfig
from tpeanuts.util.context import RuntimeContext

context = RuntimeContext.resolve("cpu", torch.float64)
oscillation = PropagationConfig.oscillation_parameters_from_preset(
    "_SM_NUFIT52_NO",
    NSI_extension="nsi_diagonal_biggio2009",   # any name in NSI_PRESETS
    context=context,
)
# oscillation.nsi is now a fully populated NSIConfig; every downstream
# builder (hamiltonian_reduced/flavour, evolutor_numerical,
# evolutor_perturbative_segment, medium.*) reads it automatically.
\end{lstlisting}
\end{examplebox}

\section{Observables and simulations that can be built with NSI}

\begin{implbox}
Every observable \texttt{tpeanuts} exposes for the Standard Model can, in
principle, be re-evaluated with \code{oscillation.nsi} set, subject to the
method-compatibility notes of \S\ref{sec:impl-pipeline}:
\begin{itemize}
  \item \textbf{Solar $P_{ee}(E)$.} \pyfunc{medium.solar.probability.solar\_probability\_state}/
    \pyfunc{solar\_probability\_mass} with \code{method="adiabatic\_exact"}
    (pointwise Hamiltonian diagonalisation, any NSI $\epsilon$, and, since
    this extension, $\epsilon^n$ whenever \code{medium.density\_n} is
    available) or \code{method="numerical"} (full coherent matrix-exponential
    propagation through the tabulated solar density profile). Used for the
    LMA/LMA-Dark degeneracy study (\cref{sec:results-lma-dark}) and the
    Borexino NSI inference fit (\cref{sec:results-borexino-nsi}).
  \item \textbf{Earth regeneration and the day--night asymmetry.}
    \pyfunc{medium.earth.probability.earth\_probability\_state}/
    \code{\_numerical}/\code{\_analytical}, propagating a solar (or any
    other) incident state through the Earth's PREM-based density profile.
    NSI in the Earth modifies the regeneration probability and hence the
    day--night asymmetry $A_{DN}$ -- the observable used to disfavour
    LMA-Dark independently of the Sun (\cref{sec:results-lma-dark}).
  \item \textbf{Atmospheric oscillograms.}
    \pyfunc{medium.atmosphere.probability.atmosphere\_probability\_transition}
    and the combined \pyfunc{pipeline.atmosphere\_earth} entry points, for a
    full atmosphere-production-to-underground-detector NSI oscillogram in
    $(E,\cos\theta_z)$.
  \item \textbf{Gradient-based inference.} \pyclass{inference.solar\_model.SolarNSIOscillationModel}
    (\cref{sec:impl-autograd}) fits a diagonal NSI coupling jointly with
    standard oscillation parameters by gradient descent against real solar
    $P_{ee}(E)$ data points, exactly as \pyclass{SolarSMOscillationModel}
    does for the pure Standard Model.
\end{itemize}
\end{implbox}

\begin{notebox}
As documented in \cref{sec:impl-pipeline}, \code{method="adiabatic\_approximated"}
(the closed-form two-level solar approximation and its Landau-Zener
correction) rejects any NSI configuration outright, and no propagation
method currently generalises the composition term to
\pyfunc{mass\_weights\_adiabatic\_approximated}.
\end{notebox}
